% !TeX root = proyecto.tex
%=========================================================
\chapter{GLOSARIO}
\begin{itemize}
    \item \textbf{API RESTful:} Interfaz de programación de aplicaciones que sigue el estilo arquitectónico REST (\textit{Representational State Transfer}), basado en convenciones para la comunicación entre cliente y servidor a través del protocolo HTTP.
    \item \textbf{bcrypt:} Función de \textit{hashing} diseñada para el almacenamiento seguro de contraseñas, que incorpora un factor de costo para dificultar ataques de fuerza bruta.
    \item \textbf{FIFO (\textit{First In, First Out}):} Criterio de priorización que establece que el primer elemento en entrar a una cola es el primero en ser atendido, respetando estrictamente el orden de llegada.
    \item \textbf{JWT (\textit{JSON Web Token}):} Estándar de autenticación sin estado que permite transmitir información verificable entre cliente y servidor mediante tokens firmados digitalmente.
    \item \textbf{KDS (\textit{Kitchen Display System}):} Interfaz de pantalla orientada al personal de cocina que presenta los pedidos pendientes, en preparación y completados, con información sobre tiempos de espera y prioridades.
    \item \textbf{KMS (\textit{Kitchen Management System}):} Sistema integral de software orientado a gestionar todos los procesos involucrados en la producción de alimentos dentro de una cocina, desde la recepción de pedidos hasta el control de inventario y el análisis del desempeño operativo.
    \item \textbf{POS (\textit{Point of Sale}):} Sistema encargado de administrar la transacción comercial en el punto de venta: el registro de productos, el cobro y la emisión de comprobantes.
    \item \textbf{Pre-ordenamiento:} Estrategia que permite a los clientes realizar su pedido con antelación, especificando un horario de recogida, con el fin de distribuir la carga de trabajo de la cocina antes del periodo de mayor afluencia.
    \item \textbf{WebSockets:} Protocolo de comunicación bidireccional y persistente que permite al servidor notificar a los clientes conectados ante cualquier cambio de estado, sin necesidad de consultas periódicas.
\end{itemize}